\section{Finite enclosures and their areas}
\label{sec:finite-footprints}

Theorem~\ref{thm:radial} supplies triangular bases and the exact limiting
cost for actual swept projections. This appendix records other finite
choices: polygonal feet with explicit area formulas, convex trapezoids,
and neighborhoods of the shortened beams' radial sectors. Each choice
keeps the same limiting cost as its sweep. Their finite areas and geometric
properties can differ, so we specify the enclosure and control its area
before taking a limit.

\subsection{One intercept identity, several enclosures}

Use a cap $R_\eps$ and fix a positive tilt for which
Appendix~\ref{sec:cap-choices} gives coverage. Retain $h=\sqrt2$.
For a closed interval $I\subset[-1,1]$ of positive length $d$, a tag $c\in I$,
and $M=\max_I R_\eps$, the upper-sweep point with parameters
$q,\rho,x$ has frozen intercept
\begin{equation}\label{eq:quadratic-foot-identity}
 u=\rho-\eps x(q-c),\qquad
 Y=qu+\frac{\eps x}{2}(q-c)^2,
 \qquad 0\le\rho\le M,\quad |x|\le2.
\end{equation}
The physical intercept point is $(0,Y,h(1-u))$. Its area element is
$h\,dY\,du$, and projection perpendicular to the frozen direction
multiplies area by $(1+\eps^2D_0(c))^{-1/2}$. The lower family is
obtained by the same orthogonal map as before.

Let $m$ be the midpoint of $I$. For positive constants $A,B$, an
absolute-value footprint is
\begin{equation}\label{eq:abs-foot-family}
 -Ad\le u\le M+Ad,\qquad
 |Y-mu|\le \frac d2|u|+Bd^2.
\end{equation}
For any tag, it contains \eqref{eq:quadratic-foot-identity} when
$A\ge2\eps$ and $B\ge\eps$. For midpoint tags, the sharper sufficient
conditions are $A\ge\eps$ and $B\ge\eps/4$.
Indeed $|q-c|\le d$ for any tag and $|q-c|\le d/2$ at the midpoint;
\eqref{eq:quadratic-foot-identity} then bounds the radial displacement
by $2\eps|q-c|$ and the extra transverse displacement by
$\eps|q-c|^2$. The angular sector itself has width $d|u|$.
The exact coordinate area is
\begin{equation}\label{eq:abs-foot-area}
 \frac{dM^2}{2}+(A+2B)Md^2+(A^2+4AB)d^3.
\end{equation}
This follows by integrating $d|u|+2Bd^2$ from $-Ad$ to $M+Ad$.
The footprint is a compact polygonal set. With a negative collar it is
generally nonconvex: the upper-boundary points at $u=\pm Ad$ have a
midpoint above the upper boundary at $u=0$.

\paragraph{The original cubic cap.}
For the cap $R_\eps=1+\eps^2H+5\eps^3$ used in
Section~\ref{sec:tetrahedron}, take $N$ equal intervals of length $d=2/N$,
midpoint $c$, and $R_I=\max_I R_\eps$. The absolute-value choice
$A=2\eps$, $B=\eps$ can be written
\begin{equation}\label{eq:polygonal-foot}
 P_I=\{(0,Y,h(1-u)):
    -2\eps d\le u\le R_I+2\eps d,\quad
    Y\in uI+[-\eps d^2,\eps d^2]\}.
\end{equation}
The intercept identity \eqref{eq:quadratic-foot-identity} proves containment.
Project $P_I$ perpendicular to $v_I=(1,-\eps G(c),h\eps c)$ and apply
$\mathcal R$ to obtain the corresponding lower-family base. These $2N$
cylinders cover the tetrahedron, and each upper base has exact area
\begin{equation}\label{eq:polygonal-area}
 \frac{h}{\sqrt{1+\eps^2D_0(c)}}
 \int_{-2\eps d}^{R_I+2\eps d}(d|u|+2\eps d^2)\,du.
\end{equation}
By \eqref{eq:abs-foot-area}, this integral is
$dR_I^2/2+O_\eps(d^2)$ uniformly over the intervals. Their total cost
therefore converges to \eqref{eq:cost-integral}. For sufficiently small
fixed $\eps$, this limit is below $2h$, so a sufficiently fine finite
partition gives a polygonal cover within the same strict area gap.
These feet are generally nonconvex;
the following affine choices give convex alternatives.

\paragraph{Affine feet.}
At a midpoint tag $c=m$ one can instead use the affine footprint
\begin{equation}\label{eq:midpoint-affine-foot}
 -Ad\le u\le M+Ad,\qquad
 |Y-cu|\le \frac d2u+Bd^2.
\end{equation}
For containment it suffices that $A\ge\eps$ and $B\ge3\eps/4$.
To check this, use the other form of the exact intercept identity,
\[
 Y-cu=(q-c)\rho-\frac{\eps x}{2}(q-c)^2.
\]
Its absolute value is at most $d\rho/2+\eps d^2/4$, and
$\rho\le u+\eps d$. These are precisely the asserted bounds.
If $B\ge A/2$, the width is nonnegative on the entire radial interval,
and the coordinate area is exactly
\begin{equation}\label{eq:affine-foot-area}
 \frac{dM^2}{2}+(A+2B)Md^2+4ABd^3.
\end{equation}
For $B>A/2$ this is a nondegenerate compact trapezoid; for
$B=A/2$ it is a nondegenerate compact triangle with its vertex at
$u=-Ad$, provided $M+2Ad>0$. All cap choices here have $M>0$.
The projection is nonsingular, so it preserves convexity and polygon type.

Freezing at a left endpoint gives another affine certificate. Put
$c=\min I$ and $W=Y-cu$. Then
\[
 W=(q-c)u+\frac{\eps x}{2}(q-c)^2.
\]
The footprint
\begin{equation}\label{eq:endpoint-affine-foot}
 -2\eps d\le u\le M+2\eps d,\qquad
 -3\eps d^2\le W\le du+3\eps d^2
\end{equation}
contains all intercepts. For $u\ge0$, the product $(q-c)u$ lies
between $0$ and $du$. For $u<0$, it lies between $du$ and $0$;
the radial bound gives $|u|d\le2\eps d^2$. The error term has
magnitude at most $\eps d^2$, proving both inequalities. This compact
convex trapezoid has area \eqref{eq:affine-foot-area} with
$A=2\eps$, $B=3\eps$.

The following choices illustrate the geometric freedom while keeping all
constants in the same physical coordinate system:
\begin{center}
\begin{tabular}{lll}
freezing tag & foot & $(A,B)$\\ \hline
arbitrary & absolute-value & $(2\eps,\eps)$\\
midpoint & absolute-value & $(\eps,\eps)$\\
midpoint & affine trapezoid & $(\eps,3\eps/4)$\\
midpoint & affine trapezoid & $(4\eps,4\eps)$\\
midpoint & affine trapezoid & $(\eps,\eps)$\\
midpoint & triangle & $(2\eps,\eps)$\\
left endpoint & affine trapezoid & $(2\eps,3\eps)$
\end{tabular}
\end{center}
For example, the midpoint trapezoid $(\eps,3\eps/4)$ is contained
in the absolute-value foot $(\eps,\eps)$ for the same cap and
interval. Their half-width difference is $\eps d^2/4$ on $u\ge0$
and $-du+\eps d^2/4$ on $u\le0$. They have the same leading
area but distinct finite areas and convexity properties.
The more generously padded choices $A=B=1$ when $\eps\le1/8$,
and $A=B=\sqrt2+2\eps$, also satisfy the absolute-value certificate.
They are useful when the original coordinate mesh is chosen independently
of the small tilt.

In all these formulas the area is $dM^2/2+O_\eps(d^2)$, uniformly
in $I$ for fixed tilt and bounded cap. Multiplication by the physical
projection factor and Riemann summation therefore give the exact total
limit $2hQ_R(\eps)$. The actual swept projections have that same limit:
the certificates give the upper bound, and the zero-free-coordinate
sectors in Theorem~\ref{thm:radial} give the matching lower bound.
This equality of limits does not identify the finite bases. For example,
freezing actual sweeps at left endpoints and projecting explicit feet at
midpoints need not give the same bases or the same finite cost.

\subsection{Two exact formulas at unit edge length}

It is useful to write two of the finite enclosures in coordinates in
which the tetrahedron has edge one. Let $h_1=1/\sqrt2$ and fix a
sufficiently small $\eta>0$. For the square-root cap
\[
 T_\eta(s)^2=\frac14+\eta^2\left(s^2/4+s^4+1/960\right),
 \qquad |s|\le\frac12,
\]
use the tetrahedron
$|x|\le(1-t)/2$, $|y|\le t/2$ in physical coordinates
$(x,y,h_1t)$. The first line family is
\[
 y=sp+\eta x(1/4+s^2),\qquad t=p+2\eta xs,
 \qquad 0\le p\le T_\eta(s),
\]
and the second exchanges $x,y$, $t,1-t$, and changes the sign of the
tilt. Under $X=4y$, $Y=4x$, $Z=2t-1$ this is cap (S) with
$\eps=\eta/2$ and $\beta=1/120$, so the coverage proof is already
complete: the first family becomes the canonical lower family with
angle $2s$ and radius $2p$; the second becomes the upper family.
For midpoint intervals of length $\Delta=1/N$, midpoint
$s_j$, and maximum $T_j$, the intercept triangles are
\begin{equation}\label{eq:unit-triangle-foot}
 -\eta\Delta\le c\le T_j+\eta\Delta,
 \qquad |b-s_jc|\le\frac\Delta2(c+\eta\Delta).
\end{equation}
Their physical area is exactly
\begin{equation}\label{eq:unit-triangle-area}
 \frac{h_1\Delta}{2}(T_j+2\eta\Delta)^2.
\end{equation}
For completeness, subtracting the frozen direction from a represented
point gives
$c=p+2\eta x(s-s_j)$ and
$b-s_jc=(s-s_j)p+\eta x(s-s_j)^2$.
Since $|x|\le1/2$,
$|c-p|\le\eta\Delta/2$ and
$|b-s_jc|\le\Delta(p+\eta\Delta/4)/2
\le\Delta(c+\eta\Delta)/2$.
Projecting these intercept triangles gives $2N$ triangular bases whose
cylinders cover for every $N\ge1$, not only after mesh refinement.
Projection divides their physical intercept area by
$\sqrt{1+\eta^2((1/4+s_j^2)^2+2s_j^2)}$.
Summing over both families gives the exact limit
\[
 h_1\int_{-1/2}^{1/2}
 \frac{1/4+\eta^2(s^2/4+s^4+1/960)}
 {\sqrt{1+\eta^2((1/4+s^2)^2+2s^2)}}\,ds
 =h_1\left(\frac14-\frac{\eta^2}{960}+O(\eta^4)\right).
\]
This is a particular triangle formula in addition to the arbitrary-tag
triangles furnished by Theorem~\ref{thm:radial}.

For the half-height construction of Appendix~\ref{sec:cap-choices}, take intervals
of half-length $\rho=1/N$, midpoint $p_j$, and maximum cap $M_j$.
Put $e=\delta\rho/2$. In its normalized intercept coordinates define
\begin{equation}\label{eq:halfheight-trapezoid}
 D_j=\{(v,w):-e\le w\le M_j+e,\quad
             |v-p_jw|\le\rho(w+2e)\}.
\end{equation}
Its exact coordinate area is
\begin{equation}\label{eq:halfheight-trapezoid-area}
 |D_j|=\rho(M_j^2+6eM_j+8e^2).
\end{equation}
An actual first-family intercept has
$(v,w)=(ur+\sigma(r+p_j)/2,u+\sigma)$ with
$\sigma=\delta X(r-p_j)/2$, so $|\sigma|\le e$ and
$|v-p_jw|\le\rho(u+e/2)\le\rho(w+2e)$.
The second-family formula has
$\sigma=-\delta Y(s-p_j)/2$ and gives the same bounds.
The perpendicular projections of these trapezoids give $2N$ compact
convex bases whose cylinders cover the tetrahedron.
The physical foot area multiplier is $h_1/2$, and its projection factor
is $(1+\delta^2(\alpha(p_j)^2+2\beta(p_j)^2))^{-1/2}$.
Consequently the cubic cap gives total area limit
\[
 \frac{h_1}{2}\int_{-1}^1
 \frac{R_\delta(r)^2}{\sqrt{1+\delta^2(\alpha(r)^2+2\beta(r)^2)}}\,dr
 =\frac{h_1}{2}\left(\frac12-\frac{\delta^2}{240}+O(\delta^3)\right).
\]
The exact cubic coefficient can also be retained by
\eqref{eq:cubic-cap-cost}. Coverage holds for $0<\delta\le1/2$;
strict saving uses a sufficiently small such $\delta$ before selecting
$N$. The count is exactly two cylinders per interval and is not asserted
to remain bounded as the tilt tends to zero.

\subsection{Beam bases from neighborhoods of two triangles}
\label{beam:neighborhood-alternative}
We finish with a different geometry for the beam bases of
Section~\ref{beam:section}. Fix $a\in(0,1/2)$ and recall
\[
 R_a=2-a,\quad c=(-1,-1),\quad r(p)=(p,1-p),\quad
 v(p)=(1,-1)+aR_a(-p^2,(1-p)^2).
\]
The physical map is $P(x,y,t)=(\beta x,\beta y,t)$ with
$\beta=1/\sqrt2$, and the first beam $H^-$ consists of
$P(c+h r(p)+t v(p),t)$ for $0\le p\le1$, $0\le h\le R_a$,
and $|t|\le1$. Its reflection in the vertical axis is the second beam.
Using the projection factor $g$ from \eqref{beam:projection-factor},
the combined limiting normalized cost is
$C_1=R_a^2\int_0^1g(v(p))\,dp$.
We construct bases from neighborhoods of two triangular sectors; their
finite areas differ from the triangular bases in the main proof.

Partition $[0,1]$ into $k$ closed intervals $I$ of length $\delta=1/k$,
and freeze the direction at each midpoint $p_I$. Sliding a point of $H^-$
to height zero along $(\beta v(p_I),1)$ leaves the planar intercept
$c+h r(p)+t(v(p)-v(p_I))$. For $e=p-p_I$, the quadratic formula for $v$
gives the exact identity
\begin{equation}\label{beam:radial-cancellation}
 v(p)-v(p_I)=-2aR_a e\,r(p)+aR_a e^2(1,-1).
\end{equation}
The first term changes only the radial parameter $h$. Indeed,
$|2aR_a et|\le\delta$ and
$|aR_a e^2t(1,-1)|\le\sqrt2\delta^2/4<\delta^2$ because
$aR_a<1$ and $|e|\le\delta/2$. Thus every intercept lies in $c+S_I$,
where $S_I$ is the closed $\delta^2$-neighborhood of
\[
 Q_I=\{d r(p):p\in I,\ -\delta\le d\le R_a+\delta\}.
\]
The set $Q_I$ consists of two triangles, with total area
\[
 |Q_I|=\frac{(R_a+\delta)^2+\delta^2}{2}\,\delta.
\]
This follows either from their vertex determinants or from the Jacobian
$|d|$ of $(d,p)\mapsto d r(p)$. Their perimeters are uniformly bounded.
For a triangle, a radius-$\rho$ neighborhood adds at most its perimeter
times $\rho$ plus $\pi\rho^2$ in area: decompose the added region
into edge rectangles and vertex sectors, whose angles sum to $2\pi$.
Applying this to both triangles,
with $\rho=\delta^2$, gives
\[
 |S_I|\le \frac{R_a^2}{2}\delta+O(\delta^2).
\]
The constant is uniform over the intervals. Projecting $P(c+S_I,0)$ onto
the perpendicular plane produces a compact base whose cylinder covers
the corresponding part of $H^-$. By \eqref{beam:projection-factor}, its normalized cost is
$|S_I|g(v(p_I))$. The unpadded sector
$\{d r(p):p\in I,\ 0\le d\le R_a\}$ lies in $Q_I\subset S_I$
and has area $R_a^2\delta/2$. Combining this lower bound with the
upper estimate gives the first beam's total normalized cost as
\[
 \frac{R_a^2}{2}\sum_I\delta\,g(v(p_I))+O_a(\delta)
 \longrightarrow \frac{C_1}{2}.
\]
Since $g\circ v$ is smooth on $[0,1]$, the midpoint sum differs from
its integral by $O_a(\delta)$. The orthogonal map
$(X,Y,T)\mapsto(-X,-Y,T)$ supplies the reflected covering at the same
cost. Thus, for every positive integer $k$, these $2k$ cylinders cover
the two beams, and their total normalized cost is
$C_1+O_a(k^{-1})$. The quadratic transverse remainder in
\eqref{beam:radial-cancellation} makes the summed padding cost vanish.
For the small tilts used in Proposition~\ref{beam:counterexample},
combining sufficiently fine beam covers of this form with the strip cover
of Lemma~\ref{beam:finite-strip} gives another finite tetrahedron cover
of total area below $A_{\min}/2$.
